\begin{table}[H]
\centering
\caption{Standardized Effect Sizes for Main Outcomes}
\label{tab:sde}
\begin{threeparttable}
\begin{adjustbox}{max width=\textwidth}
\begin{tabular}{llcccccl}
\toprule
Outcome & Specification & $\hat{\beta}$ & SE & SD($Y$) & SDE & SE(SDE) & Classification \\
\midrule
\multicolumn{8}{l}{\textit{Panel A: Pooled}} \\
Log(Planted Area) & Preferred DDD & 0.356 & 0.037 & 1.79 & 0.199 & 0.021 & Large positive \\
Log(Production) & Preferred DDD & 0.383 & 0.045 & 2.09 & 0.183 & 0.021 & Large positive \\
\midrule
\multicolumn{8}{l}{\textit{Panel B: Heterogeneous (by Pre-Reform Soybean Concentration)}} \\
Log(Planted Area) & High soy share & 0.432 & 0.038 & 1.93 & 0.224 & 0.020 & Large positive \\
Log(Planted Area) & Low soy share & 0.277 & 0.064 & 1.64 & 0.169 & 0.039 & Large positive \\
\bottomrule
\end{tabular}
\end{adjustbox}
\end{threeparttable}
\par\vspace{0.3em}
{\footnotesize \textit{Notes:} \textbf{Country:} Argentina. \textbf{Research question:} Does differential elimination of export taxes across crops cause reallocation of planted area away from the still-taxed crop toward newly liberalized crops? \textbf{Policy mechanism:} In December 2015, presidential decrees eliminated export taxes (retenciones) on wheat, corn, and sunflower while only reducing soybean export taxes by five percentage points, creating a differential incentive to shift acreage toward the newly liberalized crops. \textbf{Outcome definition:} Log of planted area in hectares at the department-crop-campaign level, from MAGyP Estimaciones Agr\'{i}colas administrative records. \textbf{Treatment:} Binary indicator equal to one for crops whose export taxes were fully eliminated (wheat, corn, sunflower) and zero for soybeans (tax reduced by five percentage points only). \textbf{Data:} MAGyP Estimaciones Agr\'{i}colas, campaigns 2010/11--2019/20, department-crop-campaign unit, 7,409 observations across 198 departments. \textbf{Method:} Difference-in-differences with department$\times$crop and department$\times$year fixed effects; standard errors clustered at department level. \textbf{Sample:} Departments growing all four focal crops (soybean, wheat, corn, sunflower) with at least four observations in both the pre-reform and post-reform periods. SDE $= \hat{\beta} / \text{SD}(Y)$ where SD($Y$) is the unconditional standard deviation of log planted area. Classification refers to magnitude, not statistical significance: Large ($|$SDE$| > 0.15$), Moderate ($0.05$--$0.15$), Small ($0.005$--$0.05$), Null ($< 0.005$).}
\end{table}

